\section{Background screening of Dirac materials}\label{app:D.41}

A Dirac target screens a dark-photon field twice, through its remote bands and through its Dirac electrons, and the scattering and absorption rates depend on where the remote bands enter. With the background in the total dielectric function the rates are smaller than in the form in which the background screens only the Dirac electrons, by $\varepsilon_{\mathrm{bg}}^2$ for scattering and $\varepsilon_{\mathrm{bg}}$ for absorption, the factors used in \cref{sec:16.1}. The notation is that of \cref{sec:7.3,app:D.22}: the half-gap $\Delta_{\mathrm{D}}$, the Dirac velocity $v_{\mathrm{D}}$ (principal velocities $v_i$), the cutoff $\Lambda_{\mathrm{D}}$, the degeneracy $g_{\mathrm{deg}}$, the independent-pair response $R^0_{\mathrm{D}}$ of \cref{eq:D.22-R} and the effective coupling $\tilde\alpha = \alpha/(\varepsilon_{\mathrm{bg}}v_{\mathrm{D}})$.

\subsection*{The total dielectric function}

The longitudinal dielectric function is $\varepsilon = 1 - v_{\mathrm{C}}\tilde R$, with $\tilde R$ the proper density response of all the charges (\cref{sec:8.4-proper}). Split $\tilde R = \tilde R_{\mathrm{D}} + \tilde R_{\mathrm{bg}}$ into the Dirac bands inside the cutoff and everything else, the remote bands and, below the optical phonons, the ionic polarization. At energies far below the gap of the remote bands and momenta small compared with the reciprocal lattice, the background responds locally and instantaneously, $1 - v_{\mathrm{C}}\tilde R_{\mathrm{bg}} \to \varepsilon_{\mathrm{bg}}$, real and independent of $\vect{q}$ and $\omega$, with local fields neglected. In the random-phase approximation the proper Dirac response is the independent-pair bubble, $\tilde R_{\mathrm{D}} = R^0_{\mathrm{D}}$, and
\begin{equation}
  \varepsilon(\vect{q},\omega) = \varepsilon_{\mathrm{bg}} + \varepsilon_{\mathrm{D}}(\vect{q},\omega), \qquad \varepsilon_{\mathrm{D}} = -v_{\mathrm{C}}(q)\,R^0_{\mathrm{D}}(\vect{q},\omega) ,
  \label{eq:D.41-eps}
\end{equation}
which is \cref{eq:7.3-epsbg}. The Dirac part $\varepsilon_{\mathrm{D}}$ does not contain $\varepsilon_{\mathrm{bg}}$: by \cref{eq:D.22-R} it is $g_{\mathrm{deg}}\alpha/(6\pi v_{\mathrm{D}})$ times the bracket of \cref{eq:D.22-eps}, and its imaginary part deep in the continuum is $g_{\mathrm{deg}}\alpha/6v_{\mathrm{D}}$. Take an anisotropic cone and a background tensor with the same principal axes. The background enters the longitudinal function as $\hat{\vect{q}}\cdot\varepsilon_{\mathrm{bg}}\cdot\hat{\vect{q}} = \sum_i\varepsilon_{\mathrm{bg},ii}\hat q_i^2$. The rescaling at the end of \cref{app:D.22} turns $q^2/v_{\mathrm{D}}$ in the overlap into $\tilde q^2/(v_1v_2v_3)$ with $\tilde{\vect{q}} = (v_1q_1, v_2q_2, v_3q_3)$, so that
\begin{equation}
  \varepsilon(\vect{q},\omega) = \sum_i\varepsilon_{\mathrm{bg},ii}\hat q_i^2 + \sum_i\hat q_i^2\,\varepsilon_{\mathrm{D},i}(\omega^2 - \tilde q^2), \qquad
  \varepsilon_{\mathrm{D},i} = \frac{g_{\mathrm{deg}}\alpha}{6\pi}\,\frac{v_i^2}{v_1v_2v_3}\,B_{\mathrm{D}}(\omega^2 - \tilde q^2) ,
  \label{eq:D.41-aniso}
\end{equation}
with $B_{\mathrm{D}}$ the bracket of \cref{eq:7.3-eps} and $v_{\mathrm{D}}\Lambda_{\mathrm{D}}$ replaced by the energy cutoff $\tilde\Lambda$ of \cref{eq:7.3-aniso}. Both terms of \cref{eq:D.41-aniso} carry the factor $\hat q_i^2$ of a longitudinal function; the background is not divided into the Dirac term.

\subsection*{Scattering}

A kinetically mixed dark photon couples to the total charge, and by \cref{eq:8.4-exact} its structure factor is $(2q^2/e^2)L/(1 - \ee^{-\beta\omega})$ with $L = \Imag(-1/\varepsilon)$ of the total dielectric function. Since $\varepsilon_{\mathrm{bg}}$ is real,
\begin{equation}
  L = \frac{\Imag\varepsilon}{|\varepsilon|^2} = \frac{\Imag\varepsilon_{\mathrm{D}}}{|\varepsilon_{\mathrm{bg}} + \varepsilon_{\mathrm{D}}|^2} .
  \label{eq:D.41-L}
\end{equation}
The numerator counts the Dirac pairs. With the spectral form of $R^0_{\mathrm{D}}$ (\cref{eq:2.4-SP}), $\Imag\varepsilon_{\mathrm{D}} = (\pi e^2/q^2)\Vol^{-1}\sum_f|\langle f|n_{\mathrm{D}}(\vect{q})|0\rangle|^2\delta(\omega - \omega_f)$ for $\omega > 0$, the golden-rule sum of pair transitions with a vertex of unit charge, and the denominator is the screening of the field by every charge in the crystal. \Cref{eq:D.41-L} is exact for the total charge within the two approximations of \cref{eq:D.41-eps}, the local background and the random-phase Dirac response, and it holds for any size of $\varepsilon_{\mathrm{D}}$.

The published Dirac rates~\cite{hochberg2018detection,geilhufe2020dirac} are assembled differently. The reference cross-section $\sigmae$ is defined in vacuum, from the free amplitude~\cite[Eqs.~(3.8), (3.10), (3.11)]{hochberg2018detection}, \cite[Eqs.~(2.3), (2.4)]{geilhufe2020dirac}. The pair transitions are summed with the same unit-charge vertex, through the overlap of the band spinors. The medium enters through the in-medium form factor $F_{\mathrm{med}} = 1/\varepsilon_{\mathrm{pub}}$, where $\varepsilon_{\mathrm{pub}} = 1 + \varepsilon_{\mathrm{D}}/\varepsilon_{\mathrm{bg}}$ is the polarization of the Dirac electrons alone, computed with the charge $e/\varepsilon_{\mathrm{bg}}^{1/2}$ in the loop and the vacuum photon propagator outside it~\cite[Eqs.~(2.6), (3.7), (3.11), (A.8)]{hochberg2018detection}, \cite[Eqs.~(2.5)--(2.7), (3.3)]{geilhufe2020dirac}. Since $\varepsilon_{\mathrm{pub}} = \varepsilon/\varepsilon_{\mathrm{bg}}$, the loss function of that construction is
\begin{equation}
  L_{\mathrm{pub}} = \frac{\Imag\varepsilon_{\mathrm{D}}}{|\varepsilon/\varepsilon_{\mathrm{bg}}|^2} = \varepsilon_{\mathrm{bg}}^2\,L ,
  \label{eq:D.41-Lpub}
\end{equation}
at every $(\vect{q},\omega)$ and for any $\varepsilon_{\mathrm{D}}$. The published scattering rate is the total-charge rate times $\varepsilon_{\mathrm{bg}}^2$ exactly.

For an anisotropic target the in-medium factor of Ref.~\cite[Eq.~(3.14)]{geilhufe2020dirac} divides each component of the Dirac polarization by its own background constant, $\varepsilon_{\mathrm{pub}} = 1 + \sum_i\hat q_i^2\varepsilon_{\mathrm{D},i}/\varepsilon_{\mathrm{bg},ii}$, and the ratio of the two loss functions is
\begin{equation}
  \frac{L_{\mathrm{pub}}}{L} = \frac{|\sum_i\hat q_i^2(\varepsilon_{\mathrm{bg},ii} + \varepsilon_{\mathrm{D},i})|^2}{|1 + \sum_i\hat q_i^2\varepsilon_{\mathrm{D},i}/\varepsilon_{\mathrm{bg},ii}|^2} .
  \label{eq:D.41-ratio}
\end{equation}
(The term of the $x$ axis is read with $q_x^2v_x^2$, where that equation prints $q_x^2v_z^2$.) Along a principal axis it is $\varepsilon_{\mathrm{bg},ii}^2$. Off the axes it is not a power of any single background constant: for $\varepsilon_{\mathrm{D}} \to 0$ it tends to $(\hat{\vect{q}}\cdot\varepsilon_{\mathrm{bg}}\cdot\hat{\vect{q}})^2$, and for finite $\varepsilon_{\mathrm{D}}$ it depends on $\hat{\vect{q}}$ through both sums. The event rates then differ by the ratio of their integrals over directions and energies, which the reach of an anisotropic target must be computed with.

\Cref{eq:D.41-Lpub} has an exact physical meaning. Let a field couple to the Dirac electrons and to no other charge. Its potential induces a Dirac charge, whose Coulomb field, screened by the background, acts back on the Dirac electrons with $v_{\mathrm{C}}/\varepsilon_{\mathrm{bg}}$; the random-phase sum of this chain is $R_{\mathrm{DD}} = R^0_{\mathrm{D}}/(1 - v_{\mathrm{C}}R^0_{\mathrm{D}}/\varepsilon_{\mathrm{bg}}) = R^0_{\mathrm{D}}/\varepsilon_{\mathrm{pub}}$. Because $R^0_{\mathrm{D}}\varepsilon_{\mathrm{pub}}^* = R^0_{\mathrm{D}} - v_{\mathrm{C}}|R^0_{\mathrm{D}}|^2/\varepsilon_{\mathrm{bg}}$ has the imaginary part $\Imag R^0_{\mathrm{D}}$, $-v_{\mathrm{C}}\Imag R_{\mathrm{DD}} = \Imag\varepsilon_{\mathrm{D}}/|\varepsilon_{\mathrm{pub}}|^2 = L_{\mathrm{pub}}$. The published rate is therefore the rate of a mediator that couples to the Dirac electrons alone. A kinetically mixed dark photon couples to the electromagnetic current, remote bands included~\cite[Eq.~(2.1)]{hochberg2018detection}, and a mediator that couples to the electron number acts on the remote-band electrons as well; with the nuclei clamped its rate is again the loss function of the total $\varepsilon$ (\cref{eq:8.4-exact}).

The factor $\varepsilon_{\mathrm{bg}}^2$ is not an artifact of the cutoff. For a gapless cone with a hard cutoff, \cref{eq:D.22-gapless} gives $\varepsilon/\varepsilon_{\mathrm{bg}} = 1 + (g_{\mathrm{deg}}\tilde\alpha/6\pi)[\ln(4v_{\mathrm{D}}^2\Lambda_{\mathrm{D}}^2/|s|) + \ii\pi\Theta(s)]$ with $s = \omega^2 - v_{\mathrm{D}}^2q^2$. The constant 1 can be absorbed into the cutoff, $\Lambda' = \Lambda_{\mathrm{D}}\exp[3\pi/(g_{\mathrm{deg}}\tilde\alpha)]$, and then $\varepsilon/\varepsilon_{\mathrm{bg}} = (g_{\mathrm{deg}}\tilde\alpha/6\pi)\mathcal{L}'$ with $\mathcal{L}' = \ln(4v_{\mathrm{D}}^2\Lambda'^2/|s|) + \ii\pi\Theta(s)$, an identity. It holds in either subtraction scheme, since the constant $\tfrac53$ by which the schemes differ is absorbed into $\Lambda'$ in the same way. It gives
\begin{equation}
  \frac{1}{|\varepsilon|^2} = \Bigl(\frac{6\pi v_{\mathrm{D}}}{g_{\mathrm{deg}}\alpha}\Bigr)^2\frac{1}{|\mathcal{L}'|^2} , \qquad
  \frac{1}{|\varepsilon_{\mathrm{pub}}|^2} = \Bigl(\frac{6\pi\varepsilon_{\mathrm{bg}}v_{\mathrm{D}}}{g_{\mathrm{deg}}\alpha}\Bigr)^2\frac{1}{|\mathcal{L}'|^2} .
  \label{eq:D.41-lprime}
\end{equation}
In the total-$\varepsilon$ rate the prefactor $g_{\mathrm{deg}}\tilde\alpha\varepsilon_{\mathrm{bg}} = g_{\mathrm{deg}}\alpha/v_{\mathrm{D}}$ does not depend on the background; in the published rate it carries $\varepsilon_{\mathrm{bg}}^2$ explicitly. The rescaled cutoff is that of Ref.~\cite[text after Eq.~(3.13)]{hochberg2018detection}, $\Lambda' = \Lambda\exp[12\pi^2\varepsilon_{\mathrm{bg}}v_{\mathrm{D}}/(g_{\mathrm{deg}}e^2)]$, and the prefactor $288\pi^4\varepsilon_{\mathrm{bg}}^2v_{\mathrm{D}}^2/(e^4g_{\mathrm{deg}}^2)$ of its Eq.~(3.13) is $\frac12(6\pi\varepsilon_{\mathrm{bg}}v_{\mathrm{D}}/g_{\mathrm{deg}}\alpha)^2$: the explicit $\varepsilon_{\mathrm{bg}}^2$ in it is the omitted $1/\varepsilon_{\mathrm{bg}}^2$. At large $\varepsilon_{\mathrm{bg}}$, $|\mathcal{L}'| \approx 6\pi/(g_{\mathrm{deg}}\tilde\alpha) \propto \varepsilon_{\mathrm{bg}}$, so the published rate tends to a constant, the unscreened rate of the Dirac pairs, while the total-$\varepsilon$ rate falls as $1/\varepsilon_{\mathrm{bg}}^2$.

\subsection*{Absorption}

A dark photon at rest, $\omega = m_\phi$ and $\vect{q} \to 0$, mixes with the photon of the medium through the self-energy of all charges, $\Pi_{\mathrm{em}} = \omega^2(1 - \varepsilon)$, the same for both polarizations at $\vect{q} \to 0$ (\cref{sec:8.4-mixing}). By \cref{eq:8.4-zetaeff}, $\zeta_{\mathrm{eff}}^2 = \zeta^2m_\phi^4/|m_\phi^2 - \Pi_{\mathrm{em}}|^2 = \zeta^2/|\varepsilon|^2$, and the absorption rate per quantum is
\begin{equation}
  \varGamma = -\frac{\zeta_{\mathrm{eff}}^2}{\omega}\,\Imag\Pi_{\mathrm{em}} = \zeta^2m_\phi\,\frac{\Imag\varepsilon}{|\varepsilon|^2} = \zeta^2m_\phi\,L(0,m_\phi) ,
  \label{eq:D.41-abs}
\end{equation}
which is \cref{eq:10.7-dp}. The published absorption rates use the self-energy of the Dirac electrons with the charge $e/\varepsilon_{\mathrm{bg}}^{1/2}$, $\Pi = \omega^2(1 - \varepsilon_{\mathrm{pub}})$~\cite[Eqs.~(4.4), (4.5)]{hochberg2018detection}, \cite[Eqs.~(2.25)--(2.27)]{geilhufe2020dirac}. Then $\zeta_{\mathrm{eff}}^2 = \zeta^2/|\varepsilon_{\mathrm{pub}}|^2$ and $-\Imag\Pi/\omega = m_\phi\Imag\varepsilon_{\mathrm{pub}} = m_\phi\Imag\varepsilon_{\mathrm{D}}/\varepsilon_{\mathrm{bg}}$, so
\begin{equation}
  \varGamma_{\mathrm{pub}} = \zeta^2m_\phi\,\frac{\Imag\varepsilon_{\mathrm{pub}}}{|\varepsilon_{\mathrm{pub}}|^2} = \zeta^2m_\phi\,\frac{\varepsilon_{\mathrm{bg}}\Imag\varepsilon_{\mathrm{D}}}{|\varepsilon|^2} = \varepsilon_{\mathrm{bg}}\,\varGamma .
  \label{eq:D.41-abspub}
\end{equation}
The power is one, not two. The absorptive part of the published self-energy already carries one factor $1/\varepsilon_{\mathrm{bg}}$, the square of the charge in the loop, whereas the scattering numerator of \cref{eq:D.41-Lpub} is the bare pair spectrum with the vacuum charge. For a dark photon polarized along the principal axis $i$ of an anisotropic target, $\varepsilon$ and $\varepsilon_{\mathrm{pub}}$ are the $ii$ components and the factor is $\varepsilon_{\mathrm{bg},ii}$. An unpolarized field averages the three rates, and the factor of the averaged rate lies between the smallest and the largest $\varepsilon_{\mathrm{bg},ii}$.

\subsection*{Reach}

In a background-free search the excluded cross-section at a given mass is inversely proportional to the expected number of events, and the excluded mixing to its square root. With \cref{eq:D.41-Lpub,eq:D.41-abspub}, the reach curves of an isotropic model correct by exact vertical shifts,
\begin{equation}
  \sigmae^{\mathrm{lim}}(m_\chi) \to \varepsilon_{\mathrm{bg}}^2\,\sigmae^{\mathrm{lim}}(m_\chi) , \qquad \zeta^{\mathrm{lim}}(m_\phi) \to \varepsilon_{\mathrm{bg}}^{1/2}\,\zeta^{\mathrm{lim}}(m_\phi) ,
  \label{eq:D.41-reach}
\end{equation}
at every mass. For an anisotropic target the shift is the ratio of the event integrals with \cref{eq:D.41-ratio} inside, which depends on the mass and on the orientation of the crystal. The published curves keep the real part of $\varepsilon_{\mathrm{D}}$ in $|\varepsilon_{\mathrm{pub}}|^2$, so \cref{eq:D.41-reach} needs no further factor. A recomputation that keeps only the leading term of \cref{eq:D.41-L} at weak screening, $L \to \Imag\varepsilon_{\mathrm{D}}/\varepsilon_{\mathrm{bg}}^2$, must restore the factor $|1 + \varepsilon_{\mathrm{D}}/\varepsilon_{\mathrm{bg}}|^{-2}$ instead, and applying both would count the screening by the Dirac electrons twice.
